\chapter{Glossary}
\label{app:glossary}
\markboth{Glossary}{}
\begin{description}[style=nextline,leftmargin=0pt,labelindent=0pt,labelwidth=0pt,labelsep=0pt,itemindent=0pt,font=\bfseries\color{Ink},itemsep=9pt]
\Needspace{5\baselineskip}\item[Admissible realization] The particular object or image allowed by a construction contract, as distinct from every compatible tuple in an ambient completion.
\Needspace{5\baselineskip}\item[Atomic diagram] The effective record of the basic relations and operations of a countable named structure. It is the standard object used to measure the degree of a copy.
\Needspace{5\baselineskip}\item[Cocycle] Here, a symmetric map $c:\Q\times\Q\to D$ satisfying the identity that makes the twisted addition on $\Q\times D$ associative.
\Needspace{5\baselineskip}\item[Cofinal quotient system] A family of finite quotients through which every finite quotient factors. Its inverse limit computes the full profinite completion.
\Needspace{5\baselineskip}\item[Computable copy] A countable realization with effective names, decidable equality, and computable operations. An oracle version permits the stated oracle.
\Needspace{5\baselineskip}\item[Computable profinite presentation] An effective inverse system of finite groups and surjective transition maps whose limit is the given profinite group. This is not a computable enumeration of all elements of an uncountable product.
\Needspace{5\baselineskip}\item[Computably enumerable set] A set whose members can be listed by a procedure. Nonmembership need not be decidable.
\Needspace{5\baselineskip}\item[Completion] An ambient object obtained by admitting compatible limits in a specified topology. In this book the relevant completion is profinite.
\Needspace{5\baselineskip}\item[Dense subgroup] A subgroup meeting every nonempty basic open set. Density does not imply equality with the ambient group.
\Needspace{5\baselineskip}\item[Divisible abelian group] An abelian group in which every equation $my=x$ with $m\ge1$ is solvable. A divisible quotient need not split off as a subgroup.
\Needspace{5\baselineskip}\item[Endomorphism-seed rank] The minimum number of elements generating a group as a module over its full endomorphism ring.
\Needspace{5\baselineskip}\item[Eventual equality] Equality modulo finitely many coordinates, written $A\eqfin B$. This is the unmarked selector invariant of $G_A$.
\Needspace{5\baselineskip}\item[Finite description] A valid finite code in a specified language with declared semantics. It need not provide computable operations.
\Needspace{5\baselineskip}\item[Finite support] A vector has only finitely many nonzero coordinates. Such vectors form $D$ inside $P$.
\Needspace{5\baselineskip}\item[Fully invariant subgroup] A subgroup preserved by every endomorphism of its ambient group. Every subgroup of the squarefree direct sum has this property.
\Needspace{5\baselineskip}\item[Grammar-relative nonrealizability] Failure to have an admissible description in one fixed grammar. It is not absolute mathematical nonexistence or indescribability.
\Needspace{5\baselineskip}\item[Marked group] A group together with named elements that isomorphisms must preserve. The mark $u_A$ retains the literal selector.
\Needspace{5\baselineskip}\item[Normal form] The unique pair $(r,d)$ representing $s_A(r)+d$. Equality of these names is simpler than the selector-dependent addition law.
\Needspace{5\baselineskip}\item[Oracle] An external set-membership resource available to a computation. Its information is charged to the effective realization contract.
\Needspace{5\baselineskip}\item[Prime spectrum] The prime orders of nonzero torsion elements. It classifies the groups $D_B$ but is constant across all $G_A$.
\Needspace{5\baselineskip}\item[Profinite completion] The inverse limit of all finite quotient groups, denoted $\widehat G$. It is not the Pontryagin dual.
\Needspace{5\baselineskip}\item[Procyclic group] A topological group with a dense cyclic subgroup. The product $P$ is procyclic.
\Needspace{5\baselineskip}\item[Pure subgroup] A subgroup $H\le M$ for which $H\cap mM=mH$ for all positive integers $m$.
\Needspace{5\baselineskip}\item[Rational line] A one-dimensional subspace over $\Q$. The preimage of $\Q q(u_A)$ is $G_A$.
\Needspace{5\baselineskip}\item[Reduced group] A group with no nonzero divisible subgroup. This property does not exclude divisible quotients.
\Needspace{5\baselineskip}\item[Residual finiteness] Every nonzero element is detected by some finite quotient. It does not imply source reconstruction from the unmarked quotient system.
\Needspace{5\baselineskip}\item[Root certificate] Either an exact element $y$ with $my=g$, or a prime divisor of $m$ at which the target has a nonzero coordinate.
\Needspace{5\baselineskip}\item[Selector] A subset $A\subseteq\N$ whose bits choose one prime from each consecutive pair. It determines the admissible rational line.
\Needspace{5\baselineskip}\item[Squarefree radical] The integer $\rad(m)$ formed by retaining each prime divisor of $m$ once.
\Needspace{5\baselineskip}\item[Tail signature] The pattern of prime-root solvability outside a finite exception set. Every infinite-order element of $G_A$ has the selector's tail signature.
\Needspace{5\baselineskip}\item[Torsion] Elements of finite order. In the product and in every $G_A$, the torsion subgroup is exactly $D$.
\Needspace{5\baselineskip}\item[Turing degree] A measure of relative computability under mutual oracle computation. The least presentation degree of $G_A$ is the degree of $A$.
\Needspace{5\baselineskip}\item[Unmarked observations] Finite quotient objects and transition maps without a named coherent source element or the original admissible image.
\end{description}
